\chapter{Introducción}\label{cap:introduccion}

\section{Contexto}
El parqueadero ``Autos Colombia'' presta su servicio exclusivamente por mensualidad y necesita un software para gestionar sus operaciones diarias: registrar las entradas y salidas de los vehículos, registrar usuarios, manejar las celdas, registrar las novedades ocurridas sobre determinado vehículo dentro del parqueadero y gestionar los pagos.

El enunciado del proyecto define tres grupos de requerimientos funcionales:
\begin{itemize}
  \item Gestión de la entrada y salida de vehículos.
  \item Gestión de usuarios y gestión de celdas.
  \item Gestión de pagos.
\end{itemize}

\section{Alcance de este documento}
El sistema se desarrolla en tres iteraciones, una por cada grupo de requerimientos. Este documento cubre la \textbf{iteración 1: gestión de la entrada y salida de vehículos}, que incluye además el registro de novedades, el control de mensualidades vencidas con el vehículo dentro y la creación de operadores, y la \textbf{iteración 2: gestión de usuarios y gestión de celdas}, que agrega propietarios y vehículos, el mapa de celdas, el cambio de propietario con historial, los horarios y turnos de los operadores y la bitácora de operaciones. En cada capítulo, el contenido de la iteración 2 aparece en una sección propia al final.

Por la decisión DEC-001, el modelo de datos se diseña completo desde la primera iteración, aunque las funcionalidades de las iteraciones 2 y 3 se implementan después.

\section{Metodología de trabajo}
El proyecto avanza de forma incremental. En cada iteración se siguen estos pasos:
\begin{enumerate}
  \item Análisis y diseño: requerimientos, historias de usuario, arquitectura, modelos de datos, diagramas UML y mockups.
  \item Implementación y pruebas del código según los artefactos aprobados.
  \item Publicación de los cambios en el repositorio de GitHub.
\end{enumerate}
Cada requerimiento, historia y decisión tiene un identificador único (RF, RNF, RN, HU, DEC) y solo se incorpora a este documento después de ser aprobado por el equipo del proyecto.

\section{Organización del documento}
El capítulo~\ref{cap:entregables} lista los entregables. El capítulo~\ref{cap:requerimientos} presenta los requerimientos y las reglas de negocio, y el capítulo~\ref{cap:arquitectura} la arquitectura del sistema. El capítulo~\ref{cap:historias} contiene las historias de usuario y el~\ref{cap:mockups} sus mockups. Los capítulos~\ref{cap:datos} a~\ref{cap:paquetes} contienen el modelo de datos y los diagramas UML. El capítulo~\ref{cap:decisiones} recoge las decisiones aprobadas y el~\ref{cap:conclusiones} las conclusiones.
